% !TEX program = xelatex
\documentclass[12pt,a4paper]{article}

% ========== بسته‌های پایه ==========
\usepackage{amsmath,amssymb,amsthm}
\usepackage{geometry}
\usepackage{enumitem}
\usepackage{xcolor}
\usepackage{graphicx}
\usepackage{setspace}
\usepackage{fancyhdr}
\usepackage{titlesec}

% ========== فونت فارسی ==========
\usepackage{fontspec}
\usepackage{xepersian}

\settextfont[Scale=1.05]{XB Yas}
\setlatintextfont[Scale=1]{Times New Roman}
\setdigitfont[Scale=1.05]{Yas}

% ========== تنظیمات صفحه ==========
\geometry{margin=2.5cm, top=3.5cm, bottom=2.5cm, headheight=25pt}
\setlength{\parindent}{0pt}
\setlength{\parskip}{0.8em}
\onehalfspacing

% ========== رنگ‌ها ==========
\definecolor{mainblue}{RGB}{25, 55, 120}
\definecolor{darkgreen}{RGB}{20, 100, 60}
\definecolor{darkred}{RGB}{140, 30, 30}

% ========== سبک عنوان‌ها ==========
\titleformat{\section}
  {\color{mainblue}\Large\bfseries}
  {\thesection}{1em}{}
  [\vspace{-0.5em}\color{mainblue}\rule{\textwidth}{1.2pt}]

\titleformat{\subsection}
  {\color{darkgreen}\large\bfseries}
  {\thesubsection}{0.8em}{}

\titleformat{\subsubsection}
  {\color{darkred}\normalsize\bfseries}
  {\thesubsubsection}{0.6em}{}

% ========== سربرگ ==========
\pagestyle{fancy}
\fancyhf{}
\fancyhead[R]{\small\color{mainblue}حل تمارین اقتصاد کلان پیشرفته}
\fancyhead[L]{\small\color{mainblue}مقطع دکتری}
\fancyfoot[C]{\thepage}
\renewcommand{\headrulewidth}{0.6pt}

% ========== دستورات ریاضی ==========
\newcommand{\R}{\mathbb{R}}
\newcommand{\MRS}{\mathrm{MRS}}

% ========== رفع تعارض tabular با xepersian ==========
\makeatletter
\let\originaltabular\tabular
\let\endoriginaltabular\endtabular
\renewenvironment{tabular}[1][c]{%
  \originaltabular[{#1}]%
}{%
  \endoriginaltabular%
}
\makeatother

% ========== شروع سند ==========
\begin{document}

\begin{center}
	{\Huge\bfseries\color{mainblue} حل کامل سؤالات ۶ و ۷}\\[0.5cm]
	{\large\color{darkgreen} با تمام گام‌های ریاضی و توضیحات مقدماتی}\\[0.3cm]
	{\color{mainblue}\rule{\textwidth}{1.5pt}}
\end{center}

\tableofcontents
\newpage

% ============================================================
\section{سؤال ۶ — اقتصاد دودوره‌ای با دو خانوار}
% ============================================================

\subsection{مقدمه: چه چیزی در مسئله داریم؟}

تصور کنید دو دوست داریم: علی (A) و بابک (B). این دو نفر در دو دوره زندگی می‌کنند: «امروز» (دوره ۱) و «فردا» (دوره ۲). هر کدام در هر دوره مقداری «کالا» دریافت می‌کنند که به آن «موهبت اولیه» یا «درآمد» می‌گوییم. کالا فسادپذیر است، یعنی اگر امروز مصرف نشود، از بین می‌رود. بنابراین برای انتقال مصرف به دوره بعد، باید از طریق «وام» یا «پس‌انداز» عمل کرد.

\textbf{درآمد علی:} در دوره ۱ مقدار ۳ و در دوره ۲ مقدار ۱.

\textbf{درآمد بابک:} در دوره ۱ مقدار ۱ و در دوره ۲ مقدار ۳.

پس علی امروز پولدار است ولی فردا فقیر، و بابک برعکس. اگر هیچ‌کس نتواند وام بگیرد یا بدهد، علی امروز زیاد مصرف می‌کند و فردا کم، و بابک برعکس. اما اگر بتوانند با هم مبادله کنند (علی امروز به بابک وام بدهد و فردا پس بگیرد)، هر دو می‌توانند مصرف خود را بین دو دوره «هموار» کنند و مطلوبیت بیشتری به‌دست آورند.

\textbf{هدف هر خانوار:} بیشینه کردن مجموع مطلوبیت مصرف امروز و فردا. مطلوبیت به‌صورت لگاریتمی است:
\[
	U_i = \ln c_{i1} + \ln c_{i2}
\]
این شکل مطلوبیت یعنی مصرف دو دوره «متقارن» هستند و خانوار دوست دارد مصرف را هموار کند. (اگر یک دوره خیلی زیاد و دوره دیگر خیلی کم باشد، مجموع لگاریتم‌ها کوچک می‌شود.)

\textbf{ابزار مبادله:} «اوراق» یا «دارایی یک‌دوره‌ای». اگر خانوار $i$ در دوره ۱ مقدار $a_i$ پس‌انداز کند (یعنی $a_i$ واحد از درآمد دوره ۱ را مصرف نکند و در قالب وام به دیگری بدهد)، در دوره ۲ مقدار $(1+r) a_i$ دریافت می‌کند. اگر $a_i < 0$ باشد، یعنی خانوار قرض گرفته است.

\textbf{قیمت وام:} نرخ بهره $r$ است. اگر $r$ مثبت باشد، پس‌اندازکننده پاداش می‌گیرد. اگر $r$ منفی باشد، پس‌اندازکننده جریمه می‌شود. تعادل اقتصاد نرخ بهره‌ای است که در آن مجموع پس‌اندازها صفر شود (چون کل اوراق عرضه خالص صفر دارد):
\[
	a_A + a_B = 0
\]

% ------------------------------------------------------------
\subsection{قسمت ۱: استخراج تقاضای دارایی هر خانوار}
% ------------------------------------------------------------

\subsubsection{گام ۱: نوشتن قید بودجه}

برای خانوار $i$ (که $i$ می‌تواند A یا B باشد):

\begin{itemize}
	\item در دوره ۱: درآمد $y_{i1}$ داریم. مقدار $a_i$ را پس‌انداز می‌کنیم. بقیه را مصرف می‌کنیم:
	      \[
		      c_{i1} = y_{i1} - a_i
	      \]
	      اگر $a_i > 0$، پس‌انداز کرده‌ایم. اگر $a_i < 0$، قرض گرفته‌ایم (یعنی $c_{i1} > y_{i1}$).

	\item در دوره ۲: درآمد $y_{i2}$ را داریم، به‌علاوه $(1+r)a_i$ که از پس‌انداز دوره ۱ به‌دست آمده:
	      \[
		      c_{i2} = y_{i2} + (1+r) a_i
	      \]
\end{itemize}

\subsubsection{گام ۲: نوشتن مسئله بهینه‌یابی}

خانوار $a_i$ را انتخاب می‌کند تا مطلوبیت کل بیشینه شود:
\[
	\max_{a_i} \ U_i = \ln(y_{i1} - a_i) + \ln(y_{i2} + (1+r) a_i)
\]

توجه کنید که $a_i$ تنها متغیر تصمیم است. $y_{i1}$، $y_{i2}$ و $r$ پارامترهای داده‌شده هستند.

\subsubsection{گام ۳: شرط مرتبه اول}

از تابع مطلوبیت نسبت به $a_i$ مشتق می‌گیریم. از قاعده زنجیره‌ای استفاده می‌کنیم:

مشتق جمله اول:
\[
	\frac{d}{da_i} \ln(y_{i1} - a_i) = \frac{1}{y_{i1} - a_i} \cdot \frac{d(y_{i1} - a_i)}{da_i} = \frac{1}{y_{i1} - a_i} \cdot (-1) = \frac{-1}{y_{i1} - a_i}
\]

مشتق جمله دوم:
\[
	\frac{d}{da_i} \ln(y_{i2} + (1+r) a_i) = \frac{1}{y_{i2} + (1+r) a_i} \cdot \frac{d(y_{i2} + (1+r)a_i)}{da_i} = \frac{1+r}{y_{i2} + (1+r)a_i}
\]

شرط مرتبه اول (مشتق کل را برابر صفر می‌گذاریم):
\[
	\frac{-1}{y_{i1} - a_i} + \frac{1+r}{y_{i2} + (1+r) a_i} = 0
\]

\subsubsection{گام ۴: حل معادله برای $a_i$ — با تمام جزئیات}

معادله را بازنویسی می‌کنیم:
\[
	\frac{1+r}{y_{i2} + (1+r) a_i} = \frac{1}{y_{i1} - a_i}
\]

دو طرف را ضربدر یکدیگر می‌کنیم (ضرب صلیبی):
\[
	(1+r)(y_{i1} - a_i) = y_{i2} + (1+r) a_i
\]

سمت چپ را باز می‌کنیم:
\[
	(1+r) y_{i1} - (1+r) a_i = y_{i2} + (1+r) a_i
\]

همه جمله‌های شامل $a_i$ را به یک طرف می‌بریم:
\[
	(1+r) y_{i1} - y_{i2} = (1+r) a_i + (1+r) a_i
\]

سمت راست را جمع می‌کنیم:
\[
	(1+r) y_{i1} - y_{i2} = 2(1+r) a_i
\]

دو طرف را بر $2(1+r)$ تقسیم می‌کنیم:
\[
	a_i = \frac{(1+r) y_{i1} - y_{i2}}{2(1+r)}
\]

این را می‌توانیم به شکل زیباتر بنویسیم:
\[
	a_i = \frac{(1+r) y_{i1}}{2(1+r)} - \frac{y_{i2}}{2(1+r)} = \frac{y_{i1}}{2} - \frac{y_{i2}}{2(1+r)}
\]

\textbf{پس فرمول نهایی تقاضای دارایی:}
\[
	\boxed{a_i = \frac{y_{i1}}{2} - \frac{y_{i2}}{2(1+r)}}
\]

\textbf{تفسیر اقتصادی:} جمله اول $\frac{y_{i1}}{2}$ یعنی خانوار دوست دارد نصف درآمد دوره ۱ را مصرف کند و نصف دیگر را پس‌انداز کند. جمله دوم $-\frac{y_{i2}}{2(1+r)}$ یعنی اگر درآمد دوره ۲ بالا باشد، خانوار امروز قرض می‌گیرد (چون می‌داند فردا می‌تواند پس بدهد). این فرمول «هموارسازی مصرف» را نشان می‌دهد.

\subsubsection{گام ۵: نوشتن تقاضای دارایی برای هر خانوار}

\textbf{برای علی (A):} $y_{A1} = 3$ و $y_{A2} = 1$
\[
	a_A = \frac{3}{2} - \frac{1}{2(1+r)}
\]

\textbf{برای بابک (B):} $y_{B1} = 1$ و $y_{B2} = 3$
\[
	a_B = \frac{1}{2} - \frac{3}{2(1+r)}
\]

\subsubsection{گام ۶: تعریف تعادل رقابتی}

تعادل رقابتی در این اقتصاد مجموعه‌ای است شامل:

\begin{itemize}
	\item یک نرخ بهره $r^*$
	\item یک تخصیص مصرف $\{c_{A1}^*, c_{A2}^*, c_{B1}^*, c_{B2}^*\}$
	\item یک تخصیص دارایی $\{a_A^*, a_B^*\}$
\end{itemize}

که در آن:

\begin{enumerate}
	\item هر خانوار با توجه به $r^*$، دارایی بهینه‌اش را انتخاب کرده باشد (یعنی فرمول $a_i$ برقرار باشد).
	\item بازار دارایی تسویه شود: $a_A^* + a_B^* = 0$.
	\item مصرف‌ها با قید بودجه سازگار باشند: $c_{i1}^* = y_{i1} - a_i^*$ و $c_{i2}^* = y_{i2} + (1+r^*) a_i^*$.
\end{enumerate}

% ------------------------------------------------------------
\subsection{قسمت ۲: تعادل بدون محدودیت استقراض}
% ------------------------------------------------------------

\subsubsection{گام ۱: اعمال شرط تسویه بازار دارایی}

\[
	a_A + a_B = 0
\]

فرمول‌های $a_A$ و $a_B$ را جایگذاری می‌کنیم:
\[
	\left[ \frac{3}{2} - \frac{1}{2(1+r)} \right] + \left[ \frac{1}{2} - \frac{3}{2(1+r)} \right] = 0
\]

\subsubsection{گام ۲: جمع کردن جمله‌های مشابه — گام به گام}

جمله‌های ثابت را جدا می‌کنیم:
\[
	\frac{3}{2} + \frac{1}{2} = \frac{4}{2} = 2
\]

جمله‌های شامل $r$ را جدا می‌کنیم:
\[
	-\frac{1}{2(1+r)} - \frac{3}{2(1+r)} = -\frac{1+3}{2(1+r)} = -\frac{4}{2(1+r)} = -\frac{2}{1+r}
\]

پس معادله تسویه به‌صورت زیر می‌شود:
\[
	2 - \frac{2}{1+r} = 0
\]

\subsubsection{گام ۳: حل برای $r$}

\[
	2 = \frac{2}{1+r}
\]

دو طرف را در $(1+r)$ ضرب می‌کنیم:
\[
	2(1+r) = 2
\]

دو طرف را بر ۲ تقسیم می‌کنیم:
\[
	1+r = 1
\]

پس:
\[
	\boxed{r = 0}
\]

\textbf{تفسیر:} نرخ بهره تعادلی صفر است. چرا؟ چون علی و بابک دقیقاً «آینه» یکدیگرند. علی امروز ۳ و فردا ۱ دارد، بابک امروز ۱ و فردا ۳ دارد. مجموع درآمد هر دو در هر دوره $3+1 = 4$ است. اگر $r=0$ باشد، علی امروز ۱ واحد به بابک وام می‌دهد (یعنی $a_A = 1$ و $a_B = -1$) و فردا همان ۱ واحد را پس می‌گیرد (بدون بهره). این باعث می‌شود مصرف هر دو در هر دوره برابر ۲ شود.

\subsubsection{گام ۴: محاسبه دارایی‌ها و مصرف‌ها}

با $r=0$:

\textbf{برای علی:}
\[
	a_A = \frac{3}{2} - \frac{1}{2(1+0)} = \frac{3}{2} - \frac{1}{2} = \frac{2}{2} = 1
\]
\[
	c_{A1} = y_{A1} - a_A = 3 - 1 = 2
\]
\[
	c_{A2} = y_{A2} + (1+r) a_A = 1 + 1 \cdot 1 = 1 + 1 = 2
\]

\textbf{برای بابک:}
\[
	a_B = \frac{1}{2} - \frac{3}{2(1+0)} = \frac{1}{2} - \frac{3}{2} = -\frac{2}{2} = -1
\]
\[
	c_{B1} = y_{B1} - a_B = 1 - (-1) = 1 + 1 = 2
\]
\[
	c_{B2} = y_{B2} + (1+r) a_B = 3 + 1 \cdot (-1) = 3 - 1 = 2
\]

\textbf{تخصیص تعادلی:}
\[
	\boxed{c_{A1} = c_{A2} = c_{B1} = c_{B2} = 2}
\]

\textbf{تفسیر:} هر دو نفر در هر دو دوره دقیقاً ۲ واحد مصرف می‌کنند. مصرف کاملاً هموار شده است. این نتیجه‌ای طبیعی است چون بازارهای مالی کامل بودند و هیچ محدودیتی وجود نداشت.

% ------------------------------------------------------------
\subsection{قسمت ۳: محدودیت استقراض برای بابک ($a_B \geq 0$)}
% ------------------------------------------------------------

\subsubsection{گام ۱: فهمیدن محدودیت}

فرض کنید بابک \textbf{نمی‌تواند قرض بگیرد}. یعنی در دوره ۱ نمی‌تواند بیش از درآمدش مصرف کند. این یعنی:
\[
	a_B \geq 0
\]
(چون اگر قرض بگیرد $a_B < 0$ می‌شود.)

علی همچنان می‌تواند آزادانه پس‌انداز کند یا قرض بگیرد.

\subsubsection{گام ۲: حدس زدن اینکه محدودیت فعال است}

اگر در تعادل بدون محدودیت، $a_B = -1$ بود (یعنی بابک قرض می‌گرفت)، حالا با محدودیت باید $a_B \geq 0$ باشد. اولین حدس این است که محدودیت \textbf{فعال} شود، یعنی $a_B = 0$.

چرا؟ چون بابک می‌خواهد قرض بگیرد (چون درآمد امروزش کم است) ولی نمی‌تواند. پس تا حد امکان، یعنی صفر، وام نمی‌گیرد.

\subsubsection{گام ۳: حل با فرض $a_B = 0$}

اگر $a_B = 0$، از تسویه بازار دارایی $a_A + a_B = 0$ نتیجه می‌گیریم:
\[
	a_A = -a_B = 0
\]

\textbf{مصرف‌ها:}

\textbf{برای بابک:}
\[
	c_{B1} = y_{B1} - a_B = 1 - 0 = 1
\]
\[
	c_{B2} = y_{B2} + (1+r) \cdot 0 = 3 + 0 = 3
\]

\textbf{برای علی:}
\[
	c_{A1} = y_{A1} - a_A = 3 - 0 = 3
\]
\[
	c_{A2} = y_{A2} + (1+r) \cdot 0 = 1 + 0 = 1
\]

\textbf{نکته مهم:} تا اینجا هیچ‌کدام مصرفشان به $r$ وابسته نیست چون $a_A = a_B = 0$ است. پس باید $r$ را از جای دیگری پیدا کنیم.

\subsubsection{گام ۴: یافتن نرخ بهره تعادلی از شرط بهینه‌یابی علی}

علی با $a_A = 0$ مجبور است مصرفش را با درآمدش برابر کند. اما آیا این بهینه است؟ شرط مرتبه اول علی می‌گوید:
\[
	\frac{1+r}{c_{A2}} = \frac{1}{c_{A1}}
\]
یعنی:
\[
	1+r = \frac{c_{A2}}{c_{A1}}
\]

با مقادیر $c_{A1} = 3$ و $c_{A2} = 1$:
\[
	1+r = \frac{1}{3}
\]
\[
	\boxed{r = \frac{1}{3} - 1 = -\frac{2}{3}}
\]

\textbf{تفسیر:} نرخ بهره منفی است! چرا؟ چون علی مجبور است امروز مصرفش را کم کند تا با $a_A = 0$ سازگار باشد. اما او ترجیح می‌دهد امروز بیشتر مصرف کند. برای اینکه علی راضی شود امروز ۳ واحد مصرف کند (به‌جای ۲ واحد در حالت بدون محدودیت)، باید نرخ بهره منفی باشد تا پس‌انداز کردن جذاب نباشد. در واقع نرخ بهره منفی باعث می‌شود علی انگیزه‌ای برای پس‌انداز نداشته باشد و مصرف امروزش بالا بماند.

\subsubsection{گام ۵: بررسی KKT برای بابک}

شرط KKT برای بابک می‌گوید:
\[
	\frac{1+r}{c_{B2}} \leq \frac{1}{c_{B1}} \quad \text{(اگر } a_B = 0 \text{)}
\]
یا به‌شکل دیگر:
\[
	1+r \leq \frac{c_{B2}}{c_{B1}}
\]

با مقادیر $c_{B1} = 1$ و $c_{B2} = 3$:
\[
	1+r \leq \frac{3}{1} = 3
\]
\[
	r \leq 2
\]

و چون $r = -2/3$ این شرط برقرار است. پس محدودیت $a_B \geq 0$ \textbf{فعال} است و $a_B = 0$ جواب بهینه است.

\subsubsection{گام ۶: بررسی دقیق‌تر — چرا $a_B = 0$ و نه چیز دیگر؟}

شرط KKT کامل برای بابک:
\[
	\frac{\partial U_B}{\partial a_B} + \mu = 0
\]
که در آن $\mu$ ضریب لاگرانژ برای قید $a_B \geq 0$ است. اگر $a_B = 0$، آنگاه $\mu > 0$ و داریم:
\[
	\frac{\partial U_B}{\partial a_B} \bigg|_{a_B=0} = \frac{-1}{c_{B1}} + \frac{1+r}{c_{B2}} = -1 + \frac{1-2/3}{3} = -1 + \frac{1/3}{3} = -1 + \frac{1}{9} = -\frac{8}{9} < 0
\]
چون مشتق منفی است، اگر بابک می‌توانست $a_B$ را کمی منفی کند (قرض بگیرد)، مطلوبیتش افزایش می‌یافت. اما نمی‌تواند! پس محدودیت فعال است و $\mu = 8/9 > 0$ است.

\textbf{تعادل جدید:}
\[
	\boxed{r = -\frac{2}{3}, \quad a_A = 0, \quad a_B = 0, \quad c_A = (3, 1), \quad c_B = (1, 3)}
\]

% ------------------------------------------------------------
\subsection{قسمت ۴: آیا تخصیص تعادل جدید پارتوکارا است؟}
% ------------------------------------------------------------

\subsubsection{گام ۱: محاسبه نرخ نهایی جانشینی}

\textbf{برای علی:}
\[
	\MRS_A = \frac{\partial U_A / \partial c_{A1}}{\partial U_A / \partial c_{A2}} = \frac{1/c_{A1}}{1/c_{A2}} = \frac{c_{A2}}{c_{A1}} = \frac{1}{3}
\]

\textbf{برای بابک:}
\[
	\MRS_B = \frac{c_{B2}}{c_{B1}} = \frac{3}{1} = 3
\]

\subsubsection{گام ۲: شرط کارایی پارتو}

تخصیص پارتوکارا است اگر و تنها اگر $\MRS_A = \MRS_B$. اما:
\[
	\frac{1}{3} \neq 3
\]

پس تخصیص \textbf{پارتوکارا نیست}.

\subsubsection{گام ۳: ساخت بهبود پارتو — گام به گام}

می‌خواهیم تخصیص جدیدی بسازیم که هم علی و هم بابک بهتر شوند. ایده این است: از علی در دوره ۱ کم کنیم و به بابک بدهیم، و از بابک در دوره ۲ کم کنیم و به علی بدهیم. این کار باعث می‌شود هر دو مصرف هموارتری داشته باشند.

\textbf{پیشنهاد:} انتقال ۱ واحد از $c_{A1}$ به $c_{B1}$، و انتقال ۱ واحد از $c_{B2}$ به $c_{A2}$.

\textbf{تخصیص جدید:}
\[
	c_A' = (3-1, 1+1) = (2, 2)
\]
\[
	c_B' = (1+1, 3-1) = (2, 2)
\]

\subsubsection{گام ۴: بررسی اینکه آیا واقعاً بهبود پارتو است}

\textbf{مطلوبیت علی قبل:}
\[
	U_A = \ln 3 + \ln 1 = \ln 3 + 0 = \ln 3 \approx 1.0986
\]

\textbf{مطلوبیت علی بعد:}
\[
	U_A' = \ln 2 + \ln 2 = 2 \ln 2 = \ln 4 \approx 1.3863
\]

چون $\ln 4 > \ln 3$، علی بهتر شده است.

\textbf{مطلوبیت بابک قبل:}
\[
	U_B = \ln 1 + \ln 3 = 0 + \ln 3 = \ln 3 \approx 1.0986
\]

\textbf{مطلوبیت بابک بعد:}
\[
	U_B' = \ln 2 + \ln 2 = 2 \ln 2 = \ln 4 \approx 1.3863
\]

بابک هم بهتر شده است.

\subsubsection{گام ۵: نتیجه}

چون هر دو بهتر شده‌اند، تخصیص تعادل جدید پارتوکارا نبود و ما یک \textbf{بهبود پارتو} ساختیم. این بهبود با انتقال ۱ واحد از هر دوره بین دو خانوار امکان‌پذیر شد.

\textbf{نکته اقتصادی:} محدودیت استقراض باعث می‌شود بابک نتواند مصرفش را هموار کند. بنابراین تخصیص ناکارا می‌شود. اگر بابک می‌توانست قرض بگیرد، مصرف هر دو برابر ۲ می‌شد و هر دو بهتر می‌شدند.

% ------------------------------------------------------------
\subsection{قسمت ۵: اگر هر دو ممنوع از استقراض باشند}
% ------------------------------------------------------------

\subsubsection{گام ۱: قیدهای جدید}

حالا هر دو خانوار نمی‌توانند قرض بگیرند:
\[
	a_A \geq 0, \quad a_B \geq 0
\]

\subsubsection{گام ۲: ترکیب با تسویه بازار}

از تسویه بازار دارایی: $a_A + a_B = 0$. اگر هر دو نامنفی باشند، تنها راه این است که هر دو صفر باشند:
\[
	a_A = 0, \quad a_B = 0
\]

چرا؟ اگر یکی مثبت باشد، دیگری باید منفی باشد تا جمع صفر شود، اما منفی ممنوع است. پس تنها راه صفر صفر است.

\subsubsection{گام ۳: مصرف‌ها}

با $a_A = a_B = 0$:
\[
	c_{A1} = 3, \quad c_{A2} = 1, \quad c_{B1} = 1, \quad c_{B2} = 3
\]

\subsubsection{گام ۴: نرخ بهره تعادلی — چرا تعیین نمی‌شود؟}

در این حالت، هیچ معامله‌ای روی نمی‌دهد. بازار دارایی وجود ندارد. بنابراین $r$ فقط از شرط بهینه‌یابی خانوارها تعیین می‌شود. اما چون هیچ‌کس معامله نمی‌کند، $r$ می‌تواند هر مقداری باشد که هر دو خانوار را راضی کند که معامله نکنند.

\textbf{شرط عدم معامله:}

علی مایل به پس‌انداز کردن نیست اگر:
\[
	\MRS_A \leq 1+r
\]
یعنی:
\[
	\frac{c_{A2}}{c_{A1}} = \frac{1}{3} \leq 1+r
\]

بابک مایل به قرض گرفتن نیست اگر:
\[
	1+r \leq \MRS_B = \frac{c_{B2}}{c_{B1}} = \frac{3}{1} = 3
\]

\textbf{بنابراین:}
\[
	\frac{1}{3} \leq 1+r \leq 3
\]

با کم کردن ۱ از همه طرف:
\[
	\boxed{-\frac{2}{3} \leq r \leq 2}
\]

\subsubsection{گام ۵: تفسیر}

هر نرخ بهره در بازه $[-2/3, 2]$ یک \textbf{تعادل} است. چرا؟ چون در این بازه، علی حاضر نیست پس‌انداز کند (چون بهره کافی برای پس‌انداز جذاب نیست) و بابک هم حاضر نیست قرض بگیرد (چون بهره برای قرض گرفتن خیلی بالاست). پس هیچ‌کس معامله نمی‌کند و بازار در تعادل است.

\textbf{نکته مهم:} در این حالت، «تعادل» فقط تخصیص مصرف را مشخص می‌کند، نه نرخ بهره را. نرخ بهره در یک بازه قرار می‌گیرد. این پدیده «عدم تعیین نرخ بهره» یا indeterminacy نام دارد.

% ============================================================
\section{سؤال ۷ — افق نامحدود}
% ============================================================

\subsection{مقدمه: چه چیزی در مسئله داریم؟}

\textbf{تصور کنید یک خانوار داریم که در زمان نامحدود زندگی می‌کند:} دوره‌های $t = 0, 1, 2, 3, \ldots$ تا بی‌نهایت. در هر دوره $t$، خانوار مقداری درآمد $y_t$ دریافت می‌کند. این درآمد می‌تواند متفاوت باشد (مثلاً در دوره ۰ مقدار ۱۰، در دوره ۱ مقدار ۵، و...). اما فرض می‌کنیم که مجموع ارزش حال درآمدها محدود است (یعنی خیلی بزرگ نیست).

\textbf{خانوار می‌تواند پس‌انداز کند یا قرض بگیرد} با نرخ بهره ثابت $R > 0$. نرخ $R$ ناخالص است (یعنی اگر ۱ واحد پس‌انداز کنید، $R$ واحد پس می‌گیرید). اگر $R > 1$، پس‌انداز سود دارد. اگر $R < 1$، پس‌انداز ضرر دارد.

\textbf{دارایی اولیه:} خانوار در ابتدای دوره ۰ مقدار $b_0$ دارایی دارد (می‌تواند مثبت یا منفی باشد).

\textbf{هدف خانوار:} بیشینه کردن مجموع مطلوبیت مصرف در همه دوره‌ها با تخفیف $\beta$:
\[
	U = \sum_{t=0}^{\infty} \beta^t \ln c_t
\]
که $0 < \beta < 1$ عامل تخفیف است. $\beta$ نزدیک به ۱ یعنی خانوار به آینده اهمیت زیادی می‌دهد. $\beta$ نزدیک به ۰ یعنی خانوار فقط به امروز اهمیت می‌دهد.

% ------------------------------------------------------------
\subsection{قسمت ۱: محدودیت ارزش حال و شرط غیرپانزی}
% ------------------------------------------------------------

\subsubsection{گام ۱: نوشتن قید بودجه دوره‌ای}

در هر دوره $t$، خانوار با محدودیت زیر روبه‌روست:
\[
	c_t + b_{t+1} = y_t + R b_t
\]

\textbf{معنای این معادله:}

\begin{itemize}
	\item سمت چپ: مصرف $c_t$ به‌علاوه دارایی که به دوره بعد می‌برد $b_{t+1}$.
	\item سمت راست: درآمد $y_t$ به‌علاوه دارایی دوره قبل با بهره $R b_t$.
\end{itemize}

به بیان دیگر، خانوار در دوره $t$ مقدار $R b_t$ از دارایی دوره قبل (اصل + بهره) دریافت می‌کند، مقدار $y_t$ درآمد دارد. این مجموع را بین مصرف $c_t$ و دارایی $b_{t+1}$ تقسیم می‌کند.

\textbf{توجه:} $b_0$ داده شده است (شرط اولیه). برای $t \geq 1$، $b_t$ متغیر تصمیم در دوره $t-1$ است.

\subsubsection{گام ۲: بازنویسی قید بودجه به‌صورت بازگشتی}

می‌خواهیم همه $b_t$ را حذف کنیم و یک محدودیت واحد بر حسب $c_t$ و $y_t$ به‌دست آوریم. برای این کار، معادله دوره $t$ را به‌صورت زیر بازنویسی می‌کنیم:
\[
	b_t = \frac{c_t + b_{t+1} - y_t}{R}
\]

این معادله بازگشتی است: $b_t$ را بر حسب $b_{t+1}$ می‌دهد. حالا این کار را برای $b_{t-1}$ و... تکرار می‌کنیم.

\subsubsection{گام ۳: جایگذاری بازگشتی — گام به گام}

\textbf{از دوره ۰ شروع می‌کنیم:}
\[
	c_0 + b_1 = y_0 + R b_0
\]
پس:
\[
	b_1 = y_0 + R b_0 - c_0
\]

\textbf{دوره ۱:}
\[
	c_1 + b_2 = y_1 + R b_1
\]
$b_1$ را از مرحله قبل جایگذاری می‌کنیم:
\[
	c_1 + b_2 = y_1 + R(y_0 + R b_0 - c_0)
\]
\[
	c_1 + b_2 = y_1 + R y_0 + R^2 b_0 - R c_0
\]
پس:
\[
	b_2 = y_1 + R y_0 + R^2 b_0 - R c_0 - c_1
\]

\textbf{دوره ۲:}
\[
	c_2 + b_3 = y_2 + R b_2
\]
$b_2$ را جایگذاری می‌کنیم:
\[
	c_2 + b_3 = y_2 + R(y_1 + R y_0 + R^2 b_0 - R c_0 - c_1)
\]
\[
	c_2 + b_3 = y_2 + R y_1 + R^2 y_0 + R^3 b_0 - R^2 c_0 - R c_1
\]
پس:
\[
	b_3 = y_2 + R y_1 + R^2 y_0 + R^3 b_0 - R^2 c_0 - R c_1 - c_2
\]

\subsubsection{گام ۴: الگو را کشف می‌کنیم}

می‌بینیم که بعد از $T$ مرحله:
\[
	b_{T+1} = \sum_{t=0}^{T} R^{T-t} y_t + R^{T+1} b_0 - \sum_{t=0}^{T} R^{T-t} c_t
\]

اگر دو طرف را بر $R^T$ تقسیم کنیم:
\[
	\frac{b_{T+1}}{R^T} = \sum_{t=0}^{T} \frac{y_t}{R^t} + R b_0 - \sum_{t=0}^{T} \frac{c_t}{R^t}
\]

یا به‌شکل دیگر:
\[
	\sum_{t=0}^{T} \frac{c_t}{R^t} + \frac{b_{T+1}}{R^T} = R b_0 + \sum_{t=0}^{T} \frac{y_t}{R^t}
\]

\subsubsection{گام ۵: حد گرفتن وقتی $T \to \infty$}

حالا $T \to \infty$ می‌بریم. اما باید مراقب جمله $\frac{b_{T+1}}{R^T}$ باشیم. این جمله می‌تواند صفر، متناهی، یا بی‌نهایت شود. برای اینکه خانوار «منطقی» باشد، باید \textbf{شرط غیرپانزی} برقرار باشد:
\[
	\lim_{T \to \infty} \frac{b_{T+1}}{R^T} \geq 0
\]

\textbf{معنای اقتصادی:} این شرط می‌گوید خانوار در نهایت نمی‌تواند بدهی بی‌پایان داشته باشد. اگر این شرط برقرار نباشد، خانوار می‌تواند با استقراض بی‌پایان، مصرف خود را بی‌نهایت کند (طرح پانزی).

\textbf{نکته ریاضی:} چرا $R^T$ در مخرج؟ چون دارایی‌های آینده با نرخ $R$ رشد می‌کنند. اگر $b_{T+1}$ با سرعت کمتر از $R$ رشد کند، نسبت به صفر میل می‌کند. اگر با سرعت $R$ رشد کند، نسبت متناهی می‌شود. اگر سریع‌تر رشد کند، نسبت به بی‌نهایت میل می‌کند که با عقلانیت سازگار نیست.

با اعمال این شرط، محدودیت ارزش حال به‌دست می‌آید:
\[
	\boxed{\sum_{t=0}^{\infty} \frac{c_t}{R^t} \leq R b_0 + \sum_{t=0}^{\infty} \frac{y_t}{R^t} = W}
\]

که $W$ \textbf{ثروت کل} خانوار است: ارزش حال کل درآمدها به‌علاوه دارایی اولیه با بهره.

% ------------------------------------------------------------
\subsection{قسمت ۲: معادله اویلر، شرط تراگذری، و مسیر مصرف}
% ------------------------------------------------------------

\subsubsection{گام ۱: نوشتن لاگرانژی}

مسئله بهینه‌یابی:
\[
	\max_{\{c_t\}_{t=0}^{\infty}} \sum_{t=0}^{\infty} \beta^t \ln c_t
	\quad \text{s.t.} \quad
	\sum_{t=0}^{\infty} \frac{c_t}{R^t} = W
\]

(قید را با تساوی می‌نویسیم چون تابع مطلوبیت اکیداً صعودی است و خانوار همه ثروتش را خرج می‌کند.)

\textbf{لاگرانژی:}
\[
	\mathcal{L} = \sum_{t=0}^{\infty} \beta^t \ln c_t + \lambda \left( W - \sum_{t=0}^{\infty} \frac{c_t}{R^t} \right)
\]

که $\lambda$ ضریب لاگرانژ است (ارزش سایه‌ای یک واحد اضافه ثروت).

\subsubsection{گام ۲: شرط مرتبه اول برای هر $c_t$}

از $\mathcal{L}$ نسبت به $c_t$ مشتق می‌گیریم. فقط دو جمله شامل $c_t$ هستند: $\beta^t \ln c_t$ و $-\lambda \frac{c_t}{R^t}$.

\[
	\frac{\partial \mathcal{L}}{\partial c_t} = \beta^t \cdot \frac{1}{c_t} - \lambda \cdot \frac{1}{R^t} = 0
\]

پس:
\[
	\frac{\beta^t}{c_t} = \frac{\lambda}{R^t}
\]

\subsubsection{گام ۳: حل برای $c_t$}

دو طرف را در $c_t$ و در $R^t$ ضرب می‌کنیم:
\[
	\beta^t R^t = \lambda c_t
\]
\[
	c_t = \frac{\beta^t R^t}{\lambda} = \frac{(\beta R)^t}{\lambda}
\]

\textbf{توجه:} $\lambda$ یک ثابت است (به $t$ بستگی ندارد). پس مصرف در دوره $t$ به‌صورت $c_t = \frac{(\beta R)^t}{\lambda}$ است.

\subsubsection{گام ۴: استخراج معادله اویلر}

نسبت مصرف دو دوره متوالی:
\[
	\frac{c_{t+1}}{c_t} = \frac{(\beta R)^{t+1}/\lambda}{(\beta R)^t/\lambda} = \frac{(\beta R)^{t+1}}{(\beta R)^t} = \beta R
\]

\textbf{پس معادله اویلر:}
\[
	\boxed{\frac{c_{t+1}}{c_t} = \beta R}
\]

\textbf{تفسیر:} اگر $\beta R > 1$، مصرف رشد می‌کند (چون صبر کردن سودآور است). اگر $\beta R < 1$، مصرف کاهش می‌یابد. اگر $\beta R = 1$، مصرف ثابت می‌ماند.

\subsubsection{گام ۵: محاسبه مسیر مصرف با استفاده از محدودیت ارزش حال}

از $c_t = (\beta R)^t c_0$ شروع می‌کنیم (چون $c_t = \frac{(\beta R)^t}{\lambda}$ و $c_0 = \frac{1}{\lambda}$، پس $c_t = (\beta R)^t c_0$).

این را در محدودیت ارزش حال جایگذاری می‌کنیم:
\[
	\sum_{t=0}^{\infty} \frac{(\beta R)^t c_0}{R^t} = W
\]

\textbf{ساده‌سازی:} $\frac{(\beta R)^t}{R^t} = \beta^t$. پس:
\[
	c_0 \sum_{t=0}^{\infty} \beta^t = W
\]

\textbf{مجموع سری هندسی:} چون $0 < \beta < 1$:
\[
	\sum_{t=0}^{\infty} \beta^t = \frac{1}{1-\beta}
\]

پس:
\[
	c_0 \cdot \frac{1}{1-\beta} = W
\]
\[
	\boxed{c_0 = (1-\beta) W}
\]

\subsubsection{گام ۶: مسیر کامل مصرف}

\[
	\boxed{c_t = (\beta R)^t c_0 = (1-\beta)(\beta R)^t W}
\]

\textbf{تفسیر:} خانوار در هر دوره بخشی از ثروت کل را مصرف می‌کند. اگر $\beta R = 1$، مصرف در همه دوره‌ها برابر $(1-\beta)W$ است (کاملاً هموار). اگر $\beta R > 1$، مصرف رشد می‌کند.

\subsubsection{گام ۷: شرط تراگذری}

\textbf{شرط تراگذری} می‌گوید که در حد، ارزش حال دارایی‌ها باید صفر شود:
\[
	\lim_{t \to \infty} \frac{b_{t+1}}{R^t} = 0
\]

این شرط از محدودیت ارزش حال با تساوی به‌دست می‌آید. اگر خانوار همه ثروتش را خرج کند، در نهایت دارایی‌اش صفر می‌شود.

\textbf{چرا شرط تراگذری مهم است؟} اگر برقرار نباشد، خانوار می‌تواند بخشی از ثروت را برای همیشه نگه دارد و مصرفش کمتر از حد بهینه باشد. شرط تراگذری تضمین می‌کند که خانوار همه ثروتش را در طول زمان خرج می‌کند.

\subsubsection{گام ۸: یکتایی پاسخ}

آیا مسیر مصرف یکتاست؟ بله. چون:

\begin{enumerate}
	\item تابع مطلوبیت اکیداً مقعر است (لگاریتمی).
	\item مجموعه بودجه محدب است.
	\item محدودیت ارزش حال با تساوی برقرار است.
\end{enumerate}

پس مسئله بهینه‌یابی یک جواب داخلی یکتا دارد.

% ------------------------------------------------------------
\subsection{قسمت ۳: مقایسه افزایش موقت و دائمی درآمد}
% ------------------------------------------------------------

\subsubsection{گام ۱: افزایش موقت در دوره $t$}

فرض کنید درآمد دوره $t$ به‌اندازه $\Delta y$ افزایش یابد (فقط یک دوره). ثروت کل $W$ به‌اندازه $\frac{\Delta y}{R^t}$ افزایش می‌یابد.

\textbf{اثر بر مصرف دوره $s$:}
\[
	\Delta c_s = (1-\beta)(\beta R)^s \cdot \frac{\Delta y}{R^t} = (1-\beta) \beta^s R^{s-t} \Delta y
\]

\textbf{اثر بر مصرف دوره $t$ (همان دوره):}
\[
	\Delta c_t = (1-\beta) \beta^t \Delta y
\]

\textbf{اثر بر مصرف دوره ۰:}
\[
	\Delta c_0 = (1-\beta) \frac{\Delta y}{R^t}
\]

\textbf{تفسیر:} افزایش موقت درآمد، مصرف را در \textbf{همه دوره‌ها} افزایش می‌دهد (به دلیل هموارسازی مصرف)، اما اثرش در دوره‌های دور کوچک‌تر است.

\subsubsection{گام ۲: افزایش دائمی در همه دوره‌ها}

فرض کنید درآمد همه دوره‌ها به‌اندازه $\Delta y$ افزایش یابد. ثروت کل $W$ به‌اندازه:
\[
	\Delta W = \sum_{t=0}^{\infty} \frac{\Delta y}{R^t} = \Delta y \cdot \frac{1}{1 - 1/R} = \Delta y \cdot \frac{R}{R-1}
\]
افزایش می‌یابد (با فرض $R > 1$).

\textbf{اثر بر مصرف دوره $s$:}
\[
	\Delta c_s = (1-\beta)(\beta R)^s \cdot \Delta y \cdot \frac{R}{R-1}
\]

\subsubsection{گام ۳: مقایسه}

\textbf{نسبت اثر دائمی به موقت (برای دوره ۰):}
\[
	\frac{\Delta c_0^{\text{دائمی}}}{\Delta c_0^{\text{موقت}}} = \frac{(1-\beta) \Delta y \cdot R/(R-1)}{(1-\beta) \Delta y / R^t} = R^t \cdot \frac{R}{R-1}
\]

برای $t=0$: نسبت $= \frac{R}{R-1}$. اگر $R=1.05$، نسبت $\approx 21$. یعنی افزایش دائمی ۱ واحد، اثرش ۲۱ برابر افزایش موقت ۱ واحد در دوره ۰ است.

\textbf{تفسیر اقتصادی:} افزایش دائمی، ثروت کل را به‌مراتب بیشتر افزایش می‌دهد و چون خانوار مصرف را هموار می‌کند، اثرش در همه دوره‌ها بزرگ‌تر است.

% ------------------------------------------------------------
\subsection{قسمت ۴: چرا معادله اویلر به‌تنهایی کافی نیست؟}
% ------------------------------------------------------------

\subsubsection{گام ۱: معادله اویلر چه می‌گوید؟}

معادله اویلر می‌گوید:
\[
	\frac{c_{t+1}}{c_t} = \beta R
\]

این یک \textbf{معادله تفاضلی} است. جواب عمومی آن:
\[
	c_t = C (\beta R)^t
\]
که $C$ یک \textbf{ثابت دلخواه} است.

\subsubsection{گام ۲: چرا $C$ دلخواه است؟}

معادله اویلر فقط نسبت مصرف دو دوره متوالی را می‌دهد، نه \textbf{سطح} مصرف. اگر $c_0 = 5$ باشد، مسیر مصرف $c_t = 5(\beta R)^t$ است. اگر $c_0 = 100$ باشد، مسیر $c_t = 100(\beta R)^t$ است. هر دو معادله اویلر را برآورده می‌کنند.

\subsubsection{گام ۳: مثال عددی برای نشان دادن مشکل}

فرض کنید $\beta = 0.9$، $R = 1.1$، و $W = 100$. پس $\beta R = 0.99 < 1$.

\textbf{اگر $C = 1$ باشد:}
\[
	c_t = (0.99)^t
\]
مجموع ارزش حال:
\[
	\sum_{t=0}^{\infty} \frac{(0.99)^t}{(1.1)^t} = \sum_{t=0}^{\infty} (0.9)^t = \frac{1}{1-0.9} = 10 \neq 100
\]
پس محدودیت بودجه برقرار نیست.

\textbf{اگر $C = 10$ باشد:}
\[
	c_t = 10(0.99)^t
\]
مجموع ارزش حال:
\[
	\sum_{t=0}^{\infty} \frac{10(0.99)^t}{(1.1)^t} = 10 \cdot 10 = 100 = W
\]
این بار محدودیت بودجه برقرار است.

\textbf{پس فقط یک مقدار $C$ (اینجا $C=10$) محدودیت بودجه را برآورده می‌کند.}

\subsubsection{گام ۴: نقش شرط تراگذری}

شرط تراگذری $\lim_{t \to \infty} \frac{b_{t+1}}{R^t} = 0$ به ما می‌گوید که خانوار نباید ثروت بی‌پایان انبار کند. این شرط همراه با محدودیت بودجه، مقدار $C$ را یکتا تعیین می‌کند.

\subsubsection{گام ۵: نتیجه}

\textbf{معادله اویلر + محدودیت بودجه + شرط تراگذری = مسیر مصرف یکتا.}

بدون محدودیت بودجه، معادله اویلر بی‌شمار جواب دارد. بدون شرط تراگذری، ممکن است جواب‌های غیربهینه (مثل انبار کردن ثروت) باقی بمانند.

\textbf{این درس مهمی است:} در مسائل بهینه‌یابی بین‌دوره‌ای، شرط مرتبه اول به‌تنهایی کافی نیست. باید شرایط مرزی (بودجه و تراگذری) را هم در نظر بگیریم.

% ------------------------------------------------------------
\subsection{خلاصه نهایی سؤال ۷}
% ------------------------------------------------------------

\begin{center}
	\begin{tabular}{|c|c|}
		\hline
		\textbf{گام}      & \textbf{نتیجه}                 \\
		\hline
		محدودیت ارزش حال  & $\sum c_t/R^t \leq W$          \\
		\hline
		معادله اویلر      & $c_{t+1}/c_t = \beta R$        \\
		\hline
		مسیر مصرف         & $c_t = (1-\beta)(\beta R)^t W$ \\
		\hline
		شرط تراگذری       & $\lim b_{t+1}/R^t = 0$         \\
		\hline
		افزایش موقت       & اثر کوچک، هموار شده            \\
		\hline
		افزایش دائمی      & اثر بزرگ، هموار شده            \\
		\hline
		ناکافی بودن اویلر & نیاز به بودجه + تراگذری        \\
		\hline
	\end{tabular}
\end{center}

\end{document}
